%%%%%%%%%%%%
%
% $Autor: Arne Bussen, Finn Wojtysiak, Francesco Kirschner, Levi Weber $
% $Datum: 2025-05 $
% $Pfad: MikrofonMP34DT05.tex $
% $Version: 4250 $
% !TeX spellcheck = en_GB/de_DE
% !TeX encoding = utf8
% !TeX root = filename 
% !TeX TXS-program:bibliography = txs:///biber
%
%%%%%%%%%%%%

% Structure
\chapter{MP34DT05} \label{MP34DT05}

In diesem Kapitel werden zunächst die Grundlagen verschiedener Mikrofontypen erläutert, um ein Verständnis für deren Funktionsweise und Einsatzmöglichkeiten zu schaffen. Im Fokus steht dabei der digitale MEMS-Mikrofonsensor MP34DT05, dessen technische Spezifikationen und Eigenschaften detailliert dargestellt werden. Ergänzend wird die Integration des Sensors mithilfe einer Implementierungsbibliothek für die Arduino IDE beschrieben, inklusive grundlegender Funktionstests und beispielhafter Anwendungscodes.

%Darüber hinaus bietet das Kapitel eine Einführung in den Aufbau und die Struktur von WAV-Dateien sowie deren Verarbeitung mit Python und MATLAB. Dabei wird gezeigt, wie sich relevante Audiofeatures aus einer WAV-Datei extrahieren, analysieren und exportieren lassen – ein zentraler Aspekt für weiterführende Anwendungen in der digitalen Signalverarbeitung.
%\Mynote{cite books, applications, board}

\section{Allgemein} \label{MP34DT05-Allgemein}

Mikrofone sind zentrale Werkzeuge der Tonaufnahme und -übertragung, deren technische Eigenschaften die Klangqualität bestimmen. Dieses Kapitel gibt einen systematischen Überblick über die wichtigsten Mikrofontypen \ref{fig:Mikrofonarten1}, ihre physikalischen Funktionsprinzipien und ihre spezifischen Einsatzgebiete.

\bigskip

Im Fokus stehen Kondensator- und dynamische Mikrofone, deren unterschiedliche Bauweisen – von Großmembran- bis Bändchenmikrofonen – klangliche und praktische Vor- und Nachteile mit sich bringen. Zudem wird die Bedeutung der Richtcharakteristik erläutert, die entscheidend beeinflusst, wie ein Mikrofon Schallquellen im Raum erfasst.

\bigskip

%\begin{figure}[ht]
%	\centering
%	\includegraphics[width=1\linewidth]{Microphon/MikrofonArten.png}
%	\caption{Liste verschiedener Mikrofonarten mit Funktionsprinzip}
%	\label{fig:Mikrofonarten}
%\end{figure}

\begin{center}
\begin{circuitikz}
    \tikzstyle{every node}=[font=\normalsize]
    
    % Dummy nodes (optional)
    \node [font=\LARGE] at (19,13.75) {};
    \node [font=\LARGE] at (17.5,17.75) {};
    \node [font=\LARGE] at (14.5,19.75) {};
    
    % Linke Kästen (dünn)
    \draw (-8.25,22.5) rectangle (-6.25,21.25);
    \draw (-8.25,20.75) rectangle (-6.25,19.75);
    \draw (-8.25,19.75) rectangle (-6.25,18.75);
    \draw (-8.25,18.25) rectangle (-6.25,17.75);
    \draw (-8.25,17.25) rectangle (-6.25,16.25);
    \draw (-8.25,14.5) rectangle (-6.25,13.5);
    
    % Restliche Kästen (dick)
    \draw [line width=1pt] (-6,22.5) rectangle (4.25,21.25);
    \draw [line width=1pt] (4.5,22.5) rectangle (16,21.25);
    \draw [line width=1pt] (-6,20.75) rectangle (-1.25,19.75);
    \draw [line width=1pt] (-6,19.75) rectangle (-3.75,18.75);
    \draw [line width=1pt] (-3.5,19.75) rectangle (-1.25,18.75);
    \draw [line width=1pt] (-1,20.75) rectangle (1,19.75);
    \draw [line width=1pt] (1.25,20.75) rectangle (4.25,19.75);
    \draw [line width=1pt] (-1,19.75) rectangle (1,18.75);
    \draw [line width=1pt] (4.5,20.75) rectangle (5.75,19.75);
    \draw [line width=1pt] (6,20.75) rectangle (8.25,19.75);
    \draw [line width=1pt] (1.25,19.75) rectangle (5.75,18.75);
    \draw [line width=1pt] (6,19.75) rectangle (8.25,18.75);
    \draw [line width=1pt] (8.5,20.75) rectangle (10,19.75);
    \draw [line width=1pt] (8.5,19.75) rectangle (10,18.75);
    \draw [line width=1pt] (10.25,20.75) rectangle (11.75,19.75);
    \draw [line width=1pt] (10.25,19.75) rectangle (11.75,18.75);
    \draw [line width=1pt] (12,20.75) rectangle (14,19.75);
    \draw [line width=1pt] (12,19.75) rectangle (14,18.75);
    \draw [line width=1pt] (14.25,20.75) rectangle (16,19.75);
    \draw [line width=1pt] (14.25,19.75) rectangle (16,18.75);
    \draw [line width=1pt] (-6,18.25) rectangle (-1.25,17.75);
    \draw [line width=1pt] (-1,18.25) rectangle (8.25,17.75);
    \draw [line width=1pt] (8.5,18.25) rectangle (14,17.75);
    \draw [line width=1pt] (14.25,18.25) rectangle (16,17.75);
    \draw [line width=1pt] (-6,17.25) rectangle (-1.25,16.25);
    \draw [line width=1pt] (-1,17.25) rectangle (11.25,16.25);
    \draw [line width=1pt] (11.5,17.25) rectangle (16,16.25);
    \draw [line width=1pt] (-3.75,14.5) rectangle (3,13.5);
    \draw [line width=1pt] (3,14.5) rectangle (9.25,13.5);
    
    % Pfeile
    \draw [->, >=Stealth] (-3.75,21.25) -- (-3.75,20.75);
    \draw [->, >=Stealth] (0,21.25) -- (0,20.75);
    \draw [->, >=Stealth] (2.75,21.25) -- (2.75,20.75);
    \draw [->, >=Stealth] (5,21.25) -- (5,20.75);
    \draw [->, >=Stealth] (7.25,21.25) -- (7.25,20.75);
    \draw [->, >=Stealth] (9.25,21.25) -- (9.25,20.75);
    \draw [->, >=Stealth] (11,21.25) -- (11,20.75);
    \draw [->, >=Stealth] (13,21.25) -- (13,20.75);
    \draw [->, >=Stealth] (15,21.25) -- (15,20.75);
    \draw [->, >=Stealth] (-5,18.75) -- (-5,18.25);
    \draw [->, >=Stealth] (-2.5,18.75) -- (-2.5,18.25);
    \draw [->, >=Stealth] (0,18.75) -- (0,18.25);
    \draw [->, >=Stealth] (3.5,18.75) -- (3.5,18.25);
    \draw [->, >=Stealth] (7.25,18.75) -- (7.25,18.25);
    \draw [->, >=Stealth] (9.25,18.75) -- (9.25,18.25);
    \draw [->, >=Stealth] (11,18.75) -- (11,18.25);
    \draw [->, >=Stealth] (13,18.75) -- (13,18.25);
    \draw [->, >=Stealth] (15,18.75) -- (15,18.25);
    \draw [->, >=Stealth] (15,17.75) -- (15,17.25);
    \draw [->, >=Stealth] (13,17.75) -- (13,17.25);
    \draw [->, >=Stealth] (10,17.75) -- (10,17.25);
    \draw [->, >=Stealth] (3.5,17.75) -- (3.5,17.25);
    \draw [->, >=Stealth] (-3.75,17.75) -- (-3.75,17.25);
    \draw [->, >=Stealth] (3,16.25) -- (3,14.5);
    \draw [->, >=Stealth] (-3.75,16.25) -- (2.5,14.75);
    \draw [->, >=Stealth] (13.75,16.25) -- (3.5,14.75);
    
    % Beschriftungen (unverändert)
    \node [font=\normalsize] at (-7.25,22.25) {\textbf{Energie-}};
    \node [font=\normalsize] at (-7.25,21.75) {\textbf{versorgung}};
    \node [font=\large] at (-0.75,21.75) {Passiv};
    \node [font=\normalsize] at (-7.25,20.5) {\textbf{Energie-}};
    \node [font=\normalsize] at (-7.5,20) {\textbf{quelle}};
    \node [font=\normalsize] at (-7.25,19.5) {\textbf{Wandler-}};
    \node [font=\normalsize] at (-7.5,19) {\textbf{prinzip}};
    \node [font=\normalsize] at (-7.25,18) {\textbf{Medium}};
    \node [font=\normalsize] at (-7.25,17) {\textbf{Arbeits-}};
    \node [font=\normalsize] at (-7.25,16.5) {\textbf{prinzip}};
    \node [font=\normalsize] at (-7.25,14.25) {\textbf{Schallfeld-}};
    \node [font=\normalsize] at (-7.5,13.75) {\textbf{größe}};
    \node [font=\normalsize] at (-3.75,20.25) {Dauermagnet};
    \node [font=\normalsize] at (-5,19.5) {Elektro-};
    \node [font=\normalsize] at (-5,19) {dynamisch};
    \node [font=\normalsize] at (-3.75,18) {Magnetfeld};
    \node [font=\normalsize] at (-3.5,17) {Geschw.-Empfänger prop.};
    \node [font=\normalsize] at (-3.5,16.5) {Membranschnelle};
    \node [font=\normalsize] at (-2.5,19.5) {Elektro-};
    \node [font=\normalsize] at (-2.5,19) {magnetisch};
    \node [font=\normalsize] at (0,20.5) {innere};
    \node [font=\normalsize] at (0,20) {Ladung};
    \node [font=\normalsize] at (0,19.5) {Piezo-};
    \node [font=\normalsize] at (0,19) {elektrisch};
    \node [font=\normalsize] at (2.75,20.5) {Dauer-};
    \node [font=\normalsize] at (2.75,20) {polarisation};
    \node [font=\normalsize] at (3.5,19.5) {Elektrostatisch};
    \node [font=\normalsize] at (3.5,19) {NF-Kondensator};
    \node [font=\normalsize] at (3.5,18) {Elektrisches Feld};
    \node [font=\normalsize] at (5,20.5) {DC-};
    \node [font=\normalsize] at (5,20) {Quelle};
    \node [font=\normalsize] at (5.25,16.75) {Auslenkungsempfänger proportional Membranauslenkung};
    \node [font=\normalsize] at (7,20.25) {AC-Quelle};
    \node [font=\normalsize] at (7.25,19.5) {Elektrostat.};
    \node [font=\normalsize] at (7.25,19) {HF-Kond.};
    \node [font=\normalsize] at (9.25,20.25) {Licht};
    \node [font=\normalsize] at (9.25,19.25) {Optisch};
    \node [font=\normalsize] at (11,20.5) {DC-};
    \node [font=\normalsize] at (11,20) {Quelle};
    \node [font=\normalsize] at (11,19.25) {Kontakt};
    \node [font=\normalsize] at (11.5,18) {Widerstand};
    \node [font=\normalsize] at (13,20.25) {Wärme};
    \node [font=\normalsize] at (13,19.5) {Thermofon};
    \node [font=\normalsize] at (13,19) {Gas};
    \node [font=\normalsize] at (13,18.75) {};
    \node [font=\normalsize] at (13.75,17) {Geschw. Empf. prop.};
    \node [font=\normalsize] at (13.75,16.5) {Membranschwelle};
    \node [font=\normalsize] at (14.25,19.75) {};
    \node [font=\normalsize] at (15,20.5) {Elektro-};
    \node [font=\normalsize] at (15,20) {magnet};
    \node [font=\normalsize] at (15,19.5) {Elektro-};
    \node [font=\normalsize] at (15,19) {magnet.};
    \node [font=\normalsize] at (15,18) {Magnet};
    \node [font=\large] at (10.5,21.75) {Aktiv};
    \node [font=\normalsize] at (-0.5,14) {Druckgradientenempfänger};
    \node [font=\normalsize] at (6.25,14) {Druckempfänger};
    
\end{circuitikz}
	\captionof{figure}{Liste verschiedener Mikrofonarten mit Funktionsprinzip}
	\label{fig:Mikrofonarten1}
\end{center}

\subsection{Ton und Frequenzen} \label{MP34DT05-TonundFrequenz}

Ton, auch bekannt als Schall, ist genauer eine einzige Schallfrequenz, die eine mechanische Deformation in dem Medium Luft verursacht, welche sich als Longitudinalwelle ausbreitet \cite{Weinzierl:2008}.

Die Frequenz $f$, gemessen in Hertz (Hz), stellt die Schwingungen über Luftdruckschwankung pro Sekunde in einer Schallwelle dar \cite{Weinzierl:2008}. Beispielsweise kann der Mensch einen Frequenzbereich von etwa 20 Hz bis 20 000 Hz wahrnehmen \cite{Weinzierl:2008}.

Es gelten folgende Formeln für die Beschreibung von Schall:

\begin{itemize}
	\item Schalldruck $p$:
	
	\begin{equation}
		p(t)=\hat{p} \cdot \cos(\omega t - \phi) 
		\quad \text{\cite{Weinzierl:2008}}
	\end{equation}
	\label{eq: Schalldruck}
	
	mit $\hat{p}$: Amplitude des Schalldrucks,
	
	$\omega = 2 \cdot \pi \cdot  f$: Kreisfrequenz, 
	
	$\phi$: Phasenverschiebung
	
	\item Schallschnelle $\nu$:
	
	\begin{equation}
		\nu(t)=\hat{\nu}\cdot \cos(\omega \cdot  t)
		\quad \text{\cite{Weinzierl:2008}}
	\end{equation}
	\label{eq: Schallschnelle}
	
	mit $\hat{\nu}$: Amplitude der Schallschnelle
	
	\item Schallgeschwindigkeit $c$ (für Gase):
	
	\begin{equation}
		c=\sqrt{\frac{\kappa p}{\rho}}
		\quad \text{\cite{Weinzierl:2008}}
	\end{equation}
	\label{eq: Schallgeschwindigkeit}
	
	\medskip
	
	Für Luft beträgt die Schallgeschwindigkeit ca. $c_{\text{Luft}} = 330\ \text{m/s}$. \cite{Bernstein:2019c}
\end{itemize}

Je schneller die Abfolge der Schwingungen, desto höher ist die Frequenz bzw. der Ton \cite{Weinzierl:2008}. Je größer die Amplitude ist, desto lauter ist der Ton \cite{Bernstein:2019b}. 

Die Lautstärke wird durch die Amplitude beschrieben, die durch den Schalldruck erzeugt wird. Die Einheit für die Lautstärke wird in Dezibel (dB) angegeben \cite{Weinzierl:2008}. Hierfür wird der Abstand der Auslenkung zum Nullpunkt gemessen. Als Nullpunkt wurde, mit $P_{\text{ref}} = 2 \cdot 10^{-5}\ \text{Pa}$ die menschliche Hörschwelle festgelegt. Ein negativer dB-Wert ist somit nicht mehr für den Menschen hörbar \cite{Weinzierl:2008}\cite{Luders:2017}.

In der Tonverarbeitung werden diese Schwingungen mechanisch, meist über Mikrofone aufgenommen \cite{Weinzierl:2008}.

\subsection{Unterschied Ton, Klang und Geräusch} \label{MP34DT05-TonKlangGeräusch}

Ein Ton im physikalischen Sinne lässt sich durch eine harmonische Schwingung mit konstanter Frequenz beschreiben. Charakteristisch ist dabei der sinusförmige Verlauf der Schalldruckänderung \cite{Weinzierl:2008}\cite{Bernstein:2019b}.

Ein Klang entsteht durch die Überlagerung eines Grundtons mit ganzzahligen Vielfachen seiner Frequenz (Obertöne). Dieses periodische Schallereignis besitzt ein diskretes Linienspektrum, wie in der harmonischen Analyse erläutert wird \cite{Weinzierl:2008}. Die spezifische Gewichtung der Obertöne bestimmt dabei die Klangfarbe.

Ein Geräusch ist durch nicht-periodische, statistisch verteilte Schwingungen gekennzeichnet \cite{Weinzierl:2008}. Im Gegensatz zu Tönen und Klängen zeigt es im Frequenzspektrum kontinuierliche Verläufe ohne ausgeprägte harmonische Struktur \cite{Bernstein:2019b}.

\bigskip

\subsection{Dezibel} \label{MP34DT05-Dezibel}
In der Audiotechnik und Elektronik wird das Dezibel (dB) zur logarithmischen Darstellung von Signalverhältnissen verwendet \cite{Weinzierl:2008}.

\subsubsection*{Grundlegende Definitionen}
Für zwei Leistungen \( P_1 \) und \( P_2 \) gilt:
\[
\text{dB} = 10 \log_{10} \left( \frac{P_2}{P_1} \right) \quad \text{\cite{Weinzierl:2008}}
\]
Bei Spannungsvergleichen (für gleiche Wellenformen und Impedanzen) ergibt sich:
\[
\text{dB} = 20 \log_{10} \left( \frac{V_2}{V_1} \right) \quad \text{\cite{Weinzierl:2008}}
\]

\subsubsection{Eigenschaften und Anwendung}
\begin{itemize}
	\item \textbf{Relatives Maß}: Dezibel beschreibt stets ein Verhältnis. Beispiele:
	\begin{itemize}
		\item \(+10\ \text{dB}\): 10-fache Leistungszunahme (\(P_2 = 10 \cdot P_1\)) \cite{Weinzierl:2008}
		\item \(-3\ \text{dB}\): Halbierung der Leistung (\(P_2 = 0{,}5 \cdot P_1\)) \cite{Weinzierl:2008}
	\end{itemize}
	
	\item \textbf{Absolutes Maß (dBm)}:
	\[
	\text{dBm} = 10 \log_{10} \left( \frac{P}{1\ \text{mW}} \right) \quad \text{\cite{Weinzierl:2008}}
	\]
	\(0\ \text{dBm}\) entspricht \(1\ \text{mW}\). In \(50\ \Omega\)-Systemen:
	\[
	V_{\text{rms}} = \sqrt{1\ \text{mW} \cdot 50\ \Omega} \approx 220\ \text{mV}_{\text{rms}} \quad \text{\cite{Weinzierl:2008}}
	\]
\end{itemize}

\subsubsection{Praxishinweise}
\begin{itemize}
	\item \textbf{Kontextabhängigkeit}: Die Interpretation von dB-Werten erfordert stets den Bezugswert (z. B. \(0\ \text{dBm} = 1\ \text{mW}\)) \cite{Weinzierl:2008}
	\item \textbf{Fehlerquellen}: Verwechslung zwischen Leistungs- (\(10 \log_{10}\)) und Spannungsverhältnissen (\(20 \log_{10}\)) vermeiden \cite{Weinzierl:2008}
\end{itemize}

\bigskip

\subsection{Richtcharakteristik von Mikrofonen}
Die Richtcharakteristik beschreibt die räumliche Empfindlichkeit eines Mikrofons gegenüber Schallwellen \cite{Weinzierl:2008}. Zu den wichtigsten Richtcharakteristiken zählen:
\begin{itemize}
	\item Kugel (omnidirektional): Empfindlichkeit für Schallwellen ist Unabhängig von der Einwirkrichtung \cite{Bernstein:2019a}.
	\item Niere (kardioid): Konzentriert sich auf frontal wirkende Schallquellen und dämpft Hintergrundgeräusche \cite{Bernstein:2019a}.
	\item Acht (bidirektional): Erfasst Schall von vorne und hinten, welcher senkrecht zur Membran einfällt. \cite{Weinzierl:2008}\cite{Bernstein:2019a}.
\end{itemize}
Der Nahbesprechungseffekt – eine Bassanhebung bei Nahaufnahme mit gerichteten Mikrofonen – wird als klangliches Gestaltungsmittel beschrieben \cite{Weinzierl:2008}.

\bigskip

\subsection{Kondensatormikrofone}

\subsubsection{Großmembranmikrofone}
Großmembranmikrofone basieren auf dem elektrostatischen Kondensatorprinzip \ref{fig:KondensatorMikrofon1}, bei dem eine dünne, leitfähige Membran und eine starre Rückplatte ein kapazitives System bilden \cite{Görne:2014d}. Dieses wird durch eine externe Spannungsversorgung (Phantompower, 48 V) aktiviert, um Schallwellen in elektrische Signale umzuwandeln \cite{Weinzierl:2008}\cite{Bernstein:2019a}. Die Membran mit einem Durchmesser von mindestens 1 Zoll (25,4 mm) erfasst mehr Schallenergie, insbesondere im Bassbereich \cite{Weinzierl:2008}.
Die hohe Empfindlichkeit und das geringe Eigenrauschen resultieren aus der effizienten Umwandlung akustischer Energie durch die große Membranfläche \cite{Weinzierl:2008}. Der lineare Frequenzgang weist eine leichte Betonung tiefer und mittlerer Frequenzen auf \cite{Weinzierl:2008}.
Ein charakteristisches Merkmal gerichteter Großmembranmikrofone ist der Proximity-Effekt, der bei Nahbesprechung zu einer Bassanhebung führt \cite{Weinzierl:2008}.
Nachteile umfassen die Anfälligkeit für mechanische Erschütterungen sowie die Abhängigkeit von Phantompower \cite{Weinzierl:2008}.

\bigskip

%\begin{figure}[ht]
%	\centering
%	\includegraphics[width=1\linewidth]{Microphon/BeispielKondensatorMikrofon.png}
%	\caption{Aufbau eines Kondensator Mikrofons}
%	\label{fig:KondensatorMikrofon}
%\end{figure}

\begin{center}
		\begin{circuitikz}
    \tikzstyle{every node}=[font=\Large]
    \node  [waves, right, rotate=0.000, line width=1pt] at(3.25,10) {};
    \draw [ color={rgb,255:red,0; green,76; blue,255}, ->, >=Stealth, dashed] (-1,8.75) -- (5.75,8.75);
    \node [font=\LARGE] at (1.25,10) {Schallwellen};
    \draw [->, >=Stealth, dashed] (21.25,8.75) -- (26,8.75);
    \draw (12.5,11.5) to[european resistor] (12.5,5.75);
    \draw (12.5,5.75) to[battery1] (7.5,5.75);
    \draw [ color={rgb,255:red,255; green,123; blue,0} ] (7.5,5.75) rectangle (7.25,10.75);
    \draw  (6.5,11.5) rectangle (6.75,5);
    \draw [short] (12.5,5.75) -- (22,5.75);
    \draw [short] (6.5,11.5) -- (22,11.5);
    \draw [short] (6.25,12) -- (20,12);
    \draw [short] (20,12) -- (20,4.5);
    \draw [short] (6.25,4.5) -- (20,4.5);
    \draw [short] (6.25,3.75) -- (20.75,3.75);
    \draw [short] (20.75,3.75) -- (20.75,12.75);
    \draw [short] (20.75,12.75) -- (6.25,12.75);
    \draw [short] (6.25,12.75) -- (6.25,12);
    \draw [short] (6.25,3.75) -- (6.25,4.5);
    \draw [->, >=Stealth] (12.5,2) -- (10.5,5.25);
    \draw [->, >=Stealth] (17.5,2) -- (13,8);
    \draw [->, >=Stealth] (10,13.25) -- (8,10.5);
    \draw [->, >=Stealth] (7.5,14.5) -- (6.75,11.75);
    \node [font=\Large] at (12.5,1.5) {Spannungsversorgung};
    \node [font=\Large] at (18.25,1.5) {hochohmiger Widerstand};
    \node [font=\Large] at (7.5,15) {Membran};
    \node [font=\Large] at (10.25,13.75) {Gegenelektrode};
    \node [font=\Large] at (23.5,9.5) {Signal};
    \draw  (-1.25,16.5) rectangle (26.25,0.75);
    \draw  (-1.25,16.5) rectangle (5,14.5);
    \node [font=\Large] at (1.75,15.75) {Beispiel Zeichnung};
    \node [font=\Large] at (2,15) {Kondensator Mikrofon:};
\end{circuitikz}

	\captionof{figure}{Aufbau eines Kondensator Mikrofons}
	\label{fig:KondensatorMikrofon1}
\end{center}

\bigskip

\subsubsection{Kleinmembranmikrofone}
Kleinmembranmikrofone nutzen ebenfalls das elektrostatische Kondensatorprinzip \ref{fig:KondensatorMikrofon1}, jedoch mit einer kleineren Membran (meist unter ¾ Zoll (19 mm) Durchmesser). Sie erfordern Phantompower (48 V) zur Aufrechterhaltung des elektrischen Feldes \cite{Weinzierl:2008}.
Durch die geringere Membranmasse erfassen sie schnelle Schalländerungen präziser \cite{Weinzierl:2008}. Sie sind weniger anfällig für den Proximity-Effekt \cite{Weinzierl:2008}.
Der lineare Frequenzgang ermöglicht eine neutrale Klangabbildung, und viele Modelle bieten umschaltbare Richtcharakteristiken \cite{Weinzierl:2008}.
Ein Nachteil ist die geringere Empfindlichkeit im Vergleich zu Großmembranmikrofonen \cite{Weinzierl:2008}.

\bigskip

\subsection{Dynamische Mikrofone}

\subsubsection{Tauchspulenmikrofon}
Tauchspulenmikrofone nutzen das elektromagnetische Induktionsgesetz: Eine Schwingspule bewegt sich im Magnetfeld eines Dauermagneten und induziert eine elektrische Spannung \cite{Weinzierl:2008}\cite{Bernstein:2019a}\cite{Görne:2014c}. Sie benötigen keine externe Spannungsversorgung \cite{Weinzierl:2008}.
Die Membran besteht aus leichten Materialien, aber die Masse der Schwingspule begrenzt den Hochtonbereich \cite{Weinzierl:2008}. Die überwiegend nierenförmige Richtcharakteristik (kardioid) unterdrückt Hintergrundgeräusche \cite{Weinzierl:2008}.
Tauchspulenmikrofone sind robust und tolerieren hohe Schalldruckpegel \cite{Weinzierl:2008}. Allerdings weisen sie eine geringere Empfindlichkeit als Kondensatormikrofone auf \cite{Weinzierl:2008}.

\bigskip

\subsubsection{Bändchenmikrofon}
Bändchenmikrofone nutzen ein ultradünnes Metallband, das im Magnetfeld eines Dauermagneten schwingt \cite{Weinzierl:2008}\cite{Bernstein:2019a}. Die bidirektionale Richtcharakteristik (Acht-Charakteristik) ist typisch \cite{Weinzierl:2008}.
Die Funktionsweise basiert auf der elektromagnetischen Induktion \cite{Görne:2014c}, wobei das Bändchen hochohmig ist und oft einen Übertrager zur Impedanzanpassung benötigt \cite{Weinzierl:2008}\cite{Bernstein:2019a}.
Bändchenmikrofone sind empfindlich gegenüber mechanischen Erschütterungen \cite{Bernstein:2019b} und hohen Schalldrücken \cite{Weinzierl:2008}.

\bigskip

\subsection{USB-Mikrofone (AD-Wandler integriert)}
USB-Mikrofone stellen eine moderne Integration von Mikrofon und Analog-Digital-Wandler in einem kompakten Gehäuse dar \cite{Weinzierl:2008}. Diese Mikrofone vereinen das elektrostatische Prinzip der Kondensatortechnik mit einem eingebauten Vorverstärker und AD-Wandler, der das im Kondensatormikrofon aufgefangene Analoge Signal in ein vom Computer lesbares Digitalsignal umwandelt \ref{fig:AD-Wandler1}, wodurch sie direkt an Computer oder mobile Geräte angeschlossen werden können \cite{Weinzierl:2008}.
Die digitale Signalverarbeitung erfolgt hier unmittelbar im Mikrofon, was eine hohe Qualität der Wandlung bei gleichzeitiger Unempfindlichkeit gegenüber elektromagnetischen Störungen auf der Übertragungsstrecke gewährleistet \cite{Weinzierl:2008}. Typischerweise arbeiten USB-Mikrofone mit Abtastraten von 44,1 kHz oder 48 kHz und einer Bittiefe von 16 oder 24 Bit \cite{Weinzierl:2008}.
Ein wesentlicher Vorteil dieser Bauform liegt in der einfachen Handhabung und Mobilität, da sie keine zusätzliche Phantomspannung oder externe Wandler benötigen \cite{Weinzierl:2008}. Allerdings sind sie in ihrer Leistung oft begrenzt und erreichen nicht die Qualität hochwertiger Studio-Kondensatormikrofone mit professionellen externen Wandlern \cite{Weinzierl:2008}.

\bigskip

%\begin{figure}[ht]
%	\centering
%	\includegraphics[width=1\linewidth]{Microphon/AD-Wandler.png}
%	\caption{Funktionsprinzip eines Analog zu Digital Wandlers}
%	\label{fig:AD-Wandler}
%\end{figure}

\begin{center}
    			\begin{circuitikz}
        \tikzstyle{every node}=[font=\LARGE]
        
        \node [font=\LARGE] at (19,13.75) {};
        \node [font=\LARGE] at (17.5,17.75) {};
        \node [font=\LARGE] at (14.5,19.75) {};
        
        \draw [ line width=1.5pt ] (-4.25,22.5) rectangle (19.25,11);
        \draw [line width=1.5pt, short] (7.25,22.5) -- (7.25,11);
        \draw [line width=1.5pt, ->, >=Stealth] (-2,13.25) -- (-2,19.5);
        \draw [line width=1.5pt, ->, >=Stealth] (-2,16.25) -- (4,16.25);
        \draw [line width=1.5pt, ->, >=Stealth] (-8.25,17.75) -- (-4.25,17.75);
        \draw [line width=1.5pt, ->, >=Stealth] (9,13.25) -- (9,19.5);
        \draw [line width=1.5pt, ->, >=Stealth] (9,16.25) -- (15.75,16.25);
        \draw [line width=1.5pt, ->, >=Stealth] (19.25,17.75) -- (23.25,17.75);
        \draw [ color={rgb,255:red,0; green,17; blue,255}, line width=1.5pt, short] (-1,16.25) .. controls (-1,19.5) and (1,19.5) .. (1,16.25);
        \draw [ color={rgb,255:red,0; green,17; blue,255}, line width=1.5pt, short] (1,16.25) .. controls (1,13.25) and (3,13.25) .. (3,16.25);
        \draw [ color={rgb,255:red,112; green,112; blue,112}, line width=1pt, dashed] (10,16.25) .. controls (9.75,19.75) and (12.25,19.75) .. (12,16.25);
        \draw [ color={rgb,255:red,112; green,112; blue,112}, line width=1pt, dashed] (12,16.25) .. controls (11.75,12.75) and (14.25,12.75) .. (14,16.25);
        \draw [ color={rgb,255:red,0; green,17; blue,255} , line width=1pt ] (10,16.25) rectangle (10.5,17.25);
        \draw [ color={rgb,255:red,0; green,17; blue,255} , line width=1pt ] (10.5,16.25) rectangle (11,18.5);
        \draw [ color={rgb,255:red,0; green,17; blue,255} , line width=1pt ] (11,16.25) rectangle (11.5,18.5);
        \draw [ color={rgb,255:red,0; green,17; blue,255} , line width=1pt ] (11.5,16.25) rectangle (12,17.25);
        \draw [ color={rgb,255:red,0; green,17; blue,255} , line width=1pt ] (12,15.25) rectangle (12.5,16.25);
        \draw [ color={rgb,255:red,0; green,17; blue,255} , line width=1pt ] (13.5,15.25) rectangle (14,16.25);
        \draw [ color={rgb,255:red,0; green,17; blue,255} , line width=1pt ] (12.5,14) rectangle (13,16.25);
        \draw [ color={rgb,255:red,0; green,17; blue,255} , line width=1pt ] (13,14) rectangle (13.5,16.25);
        \node [font=\LARGE] at (7.25,23.25) {AD-Wandler};
        \node [font=\LARGE] at (-6.25,18.25) {Eingang};
        \node [font=\LARGE] at (-2.5,19.5) {U};
        \node [font=\LARGE] at (0.5,19.5) {Analogsignal:};
        \node [font=\LARGE] at (4.5,16.25) {t};
        \node [font=\LARGE] at (8.5,19.5) {U};
        \node [font=\LARGE] at (12,19.5) {Digitalsignal:};
        \node [font=\LARGE] at (16.25,16.25) {t};
        \node [font=\LARGE] at (21.25,18.25) {Ausgang};
    \end{circuitikz}
    
	\captionof{figure}{Funktionsprinzip eines Analog zu Digital Wandlers}
	\label{fig:AD-Wandler1}
\end{center}

\bigskip

\subsection{Mikrofone mit externen AD-Wandlern}
Für hochqualitative digitale Aufnahmen kommen zunehmend Mikrofone mit externen AD-Wandlern zum Einsatz \cite{Weinzierl:2008}. Diese ermöglichen durch getrennte hochwertige Analog-Digital-Wandlung eine rauscharme und präzise Signalübertragung, insbesondere bei langen Kabelstrecken \cite{Weinzierl:2008}. 

\bigskip

\subsection{MEMS-Mikrofon}
Ein MEMS-Mikrofon (Micro-Electro-Mechanical Systems) ist ein miniaturisiertes Mikrofon, das mechanische und elektronische Komponenten auf einem Siliziumchip vereint. Es wandelt Schallwellen über eine bewegliche Membran und eine feste Gegenelektrode in elektrische Signale um, wobei die durch Schall verursachten Abstandsänderungen kapazitiv gemessen werden \cite{Weinzierl:2008}. Moderne MEMS-Mikrofone liefern digitale Audiosignale und nutzen dabei häufig Modulationstechniken wie PDM (Pulse Density Modulation) und anschließend PCM (Pulse Code Modulation) im angeschlossenen Mikrocontroler , um die analogen Schwingungen in ein digitales Format zu übertragen. Diese Methoden ermöglichen eine effiziente Integration in digitale Systeme, etwa in Smartphones oder IoT-Geräte, und tragen zur kompakten Bauweise sowie geringem Stromverbrauch bei \cite{Weinzierl:2008}.

\bigskip

\subsection{Pulsdichtenmodulation (PDM)}
Die Pulsdichtenmodulation (PDM) ist ein digitales Modulationsverfahren zur Umwandlung analoger Signale in einen binären Bitstrom \ref{fig:PDM}. Im Gegensatz zur Puls-Code-Modulation (PCM) verwendet PDM eine 1-Bit-Quantisierung bei extrem hohen Abtastraten (typischerweise im MHz-Bereich). Der analoge Signalwert wird durch die zeitliche Dichte logischer „1"-Pulse kodiert: Ein höherer Momentanwert des Eingangssignals führt zu einer dichteren Abfolge von „1"-Signalen, während niedrigere Amplituden vermehrt „0"-Anteile enthalten \cite{Weinzierl:2008}. Im folgenden wird auf die PDM Methode genauer eingegangen.

\begin{center}
    \usetikzlibrary{arrows.meta, decorations.pathmorphing}
    
    
    \centering
    \resizebox{1\textwidth}{!}{%
        \begin{circuitikz}
            \tikzstyle{every node}=[font=\Huge]
            \draw  (20.5,7.75) rectangle node {\LARGE  } (20.5,7.75);
            \node [font=\LARGE] at (20.5,7.75) {};
            \draw  (1.25,9.75) rectangle (6.25,6.75);
            \draw [->, >=Stealth] (6.5,8.25) -- (8.75,8.25);
            \draw  (9.25,9.75) rectangle (14,6.75);
            \node [font=\LARGE] at (3.75,8.75) {Kondensator-};
            %\draw [decorate, decoration={snake, amplitude=2pt, segment length=3pt}, line width=1pt] (-2.25,8.25) -- (-0.5,8.25);
            \node [font=\LARGE] at (3.75,8) {mikrofon};
            \draw [line width=0.5pt, short] (7.5,8.25) -- (7.5,11.5);
            \node [font=\LARGE] at (-1,9.5) {Audiowellen};
            \node  [waves, right, rotate=0.000, line width=1pt] at(-2.25,8.25) {};
            \draw [line width=0.5pt, ->, >=Stealth] (-0.25,8.25) -- (1,8.25);
            \draw [ line width=0.5pt ] (4.75,13.25) rectangle (10.25,11.5);
            \node [font=\Huge] at (7.5,12.25) {Analoges Signal};
            \node [font=\LARGE] at (10,8.75) {};
            \draw[domain=5.5:10,samples=100,smooth] plot (\x,{1*sin(10*\x r -5.4 r ) +14.75});
            \node [font=\Huge] at (11.5,8.25) {PDM};
            \draw [->, >=Stealth] (14.25,8.25) -- (16.75,8.25);
            \draw  (17,6.75) rectangle (23.5,9.75);
            \foreach \x in {0,...,5}{
                \draw  (13.5+\x*0.7142857142857143,14) -- ++(0,1.7) -- ++ (0.35714285714285715, 0) -- ++(0, -1.7) -- ++(0.35714285714285715,0);
            }
            \draw [short] (15.5,8.25) -- (15.5,11.5);
            \draw  (12.5,13.25) rectangle (18.5,11.5);
            \node [font=\Huge] at (15.5,12.25) {Digitales Signal};
            \node [font=\Huge] at (20.25,8.25) {Weiterverarbeitung};
        \end{circuitikz}
    }%
    \label{fig:my_lable}
    
	\captionof{figure}{Umwandlung eines Signals mittels PDM}
	\label{fig:PDM}
\end{center}

\subsubsection{Funktionsprinzip}
PDM basiert auf drei Schlüsselelementen:
\begin{itemize}
	\item \textbf{Oversampling}: Die Abtastrate $f_s$ liegt weit über dem Nyquist-Kriterium (Abtastrate mindestens doppelt so hoch wie die Signalfrequenz) (\cite{Görne:2014a})(z. B. dem 64-fachen der Zielbandbreite). Dies verschiebt das Quantisierungsrauschen in hochfrequente Bereiche außerhalb des Nutzbandes \cite{Weinzierl:2008}\cite{Görne:2014b}.
	
	\item \textbf{Rauschformung (Noise Shaping)}: Durch eine Rückkopplungsschleife im Delta-Sigma-Modulator wird das Quantisierungsrauschen spektral umgeformt. Tiefe Frequenzen im Nutzsignalbereich bleiben rauscharm, während hochfrequentes Rauschen dominiert und später durch digitale Filter entfernt wird \cite{Weinzierl:2008}.
	
	\item \textbf{1-Bit-Quantisierung}: Ein Komparator wandelt das analoge Signal in einen Pulsstrom um. Dies reduziert die Hardwarekomplexität im Vergleich zu mehrbitigen Analog-Digital-Wandlern (ADCs) erheblich \cite{Weinzierl:2008}.
\end{itemize}

\subsubsection{Mathematische Grundlagen}
Die PDM lässt sich durch folgende Gleichungen beschreiben:
\begin{itemize}
	\item Übertragungsfunktion des Delta-Sigma-Modulators:
	\begin{equation}
		Y(z) = X(z) \cdot z^{-1} + E(z) \cdot \left(1 - z^{-1}\right),
	\end{equation}
	wobei $Y(z)$ den Ausgangsbitstrom, $X(z)$ das Eingangssignal und $E(z)$ das Quantisierungsrauschen darstellt \cite{Weinzierl:2008}.
	
	\item Das Signal-Rausch-Verhältnis (SNR) im Basisband hängt von der Oversampling-Rate $N$ und der Ordnung $L$ des Rauschformers ab:
	\begin{equation}
		\text{SNR} \propto 20 \cdot \log_{10}\left(N^{L+0.5}\right) \, \text{dB}.
	\end{equation}
	\cite{Weinzierl:2008}.
\end{itemize}

\subsubsection{Anwendungen}
PDM wird primär in Systemen mit begrenzten Hardware-Ressourcen eingesetzt, z. B. in digitalen Audiosensoren oder integrierten Schaltkreisen. Die Technik ist ideal für Anwendungen, die geringen Stromverbrauch und minimale Pin-Belegung erfordern \cite{Weinzierl:2008}.

\bigskip

\subsection{Puls-Code-Modulation (PCM)}
Die \textbf{Puls-Code-Modulation (PCM)} ist ein lineares Digitalisierungsverfahren, das analoge Audiosignale durch Abtastung, Quantisierung und Binärkodierung in ein digitales Format überführt. Dies passiert jedoch nicht im eigentlichen Mikrofon an sich, sondern im daran angeschlosssenem Mikrocontroller, beispielsweise Arduino, der für eine weitere Verarbeitung des Signals einen Code und nicht ein Digitalsignal benötigt. Es gilt als Standard in der professionellen Audiotechnik, da es eine verlustfreie Signalabbildung ermöglicht \cite{Weinzierl:2008}.

\subsubsection{Funktionsprinzip}
PCM basiert auf drei zentralen Schritten:
\begin{itemize}
	\item \textbf{Abtastung (Sampling)}: Das analoge Signal wird mit einer festen Abtastrate \( f_s \) (z.\,B. 44,1\,kHz oder 48\,kHz) diskretisiert. Gemäß dem \textbf{Nyquist-Shannon-Theorem} muss \( f_s \geq 2 \cdot f_{\text{max}} \) gelten, um Aliasing zu vermeiden \cite{Weinzierl:2008}.
	
	\item \textbf{Quantisierung}: Jeder Abtastwert wird auf einen diskreten Wert aus \( 2^b \) möglichen Stufen abgebildet, wobei \( b \) die Bittiefe (z.\,B. 16\,Bit, 24\,Bit) definiert. Höhere Bittiefen reduzieren das Quantisierungsrauschen und erhöhen den Dynamikumfang \cite{Weinzierl:2008}.
	
	\item \textbf{Kodierung}: Die quantisierten Werte werden als Binärzahlen (Pulse) gespeichert oder übertragen. PCM nutzt eine \textbf{lineare Quantisierungskennlinie}, bei der alle Amplitudenstufen gleichabständig sind \cite{Weinzierl:2008}.
\end{itemize}

\subsubsection{Mathematische Grundlagen}
Es lassen sich folgende Kernparameter definieren:
\begin{itemize}
	\item \textbf{Bitrate \( R \)}:
	\begin{equation}
		R = f_s \cdot b \cdot k \quad \text{[Bit/s]},
	\end{equation}
	wobei \( k \) die Anzahl der Audiokanäle ist \cite{Weinzierl:2008}.
	
	\item \textbf{Signal-Rausch-Verhältnis (SNR)}: Bei linearer Quantisierung gilt:
	\begin{equation}
		\text{SNR} \approx 6.02 \cdot b + 1.76 \quad \text{dB}.
	\end{equation}
	Ein 16-Bit-System erreicht somit ca. 98\,dB SNR \cite{Weinzierl:2008}.
	
	\item \textbf{Dynamikumfang}: Theoretisch \( \approx 6 \cdot b \) dB, da jede Verdopplung der Amplitude 6\,dB entspricht \cite{Weinzierl:2008}.
\end{itemize}

\subsubsection{Anwendungen}
PCM ist die Basis für:
\begin{itemize}
	\item \textbf{Audio-CDs} (44,1\,kHz, 16\,Bit),
	\item \textbf{Digitale Tonstudios} (DAWs, 96\,kHz, 24\,Bit),
	\item \textbf{Telekommunikation} (VoIP, ISDN),
	\item \textbf{High-Fidelity-Audiogeräte} (DACs, digitale Mixer).
\end{itemize}
Aufgrund seiner universellen Kompatibilität bleibt PCM trotz effizienterer Methoden wie PDM in präzisionskritischen Anwendungen dominant \cite{Weinzierl:2008}.

\bigskip

\section{Spezifischer Sensor MP34DT05}
Das MP34DT05 ist ein kompaktes, digitales Mikrofon mit omnidirektionaler Schallaufnahme, welches für moderne Audioanwendungen entwickelt wurde und als integrierter Sensor des Arduino Nano 33 BLE Sense Lite bereits verbaut ist \ref{fig:microphone_detail}. Es kombiniert einen kapazitiven MEMS-Sensor mit einer integrierten CMOS-Schnittstelle und liefert ein digitales PDM-Ausgangssignal (Pulse Density Modulation). Das Mikrofon hat eine geringe Baugröße (HCLGA-Gehäuse: 4x3x1 mm), einen geringen Stromverbrauch und ist robust, weshalb es ideal für platzbeschränkte und mobile Geräte ist. (Datenblatt MP34DT05, S. 1).

\begin{figure}[h]
	\centering
	
\begin{tikzpicture}
		\ArduinoNanoTikz;

		\fill[gray, opacity=0.7] (-11.2,-0.2) rectangle (0.1,4.7);

		\coordinate (A) at (-4.3,2.55);
		\coordinate (B) at (-3.3,1.85);    
		
		\begin{scope}
			\clip (A) rectangle (B);
			\ArduinoNanoTikz
		\end{scope}
		
	\end{tikzpicture}
	\caption{MP34DT05-Mikrofon auf dem Arduino Nano 33 BLE Sense}
	\label{fig:microphone_detail}
\end{figure}

\bigskip

\subsection{Akustische Leistung}
Das MP34DT05 deckt einen Frequenzbereich von 100 Hz bis 10 kHz ab, wobei die typische Frequenzantwort auf 1KHz normalisiert ist (Datenblatt MP34DT05, S. 8, Abb. 3. Bei einem Schalldruckpegel von 94 dB SPL erreicht das MP34DT05 eine Empfindlichkeit von -25 dBFS ± 3 dB und einen Signal-Rausch-Abstand (SNR) von 64 dB. Der maximale Schalldruck liegt bei 122,5 dB SPL, wodurch auch laute Umgebungen erfasst werden können \cite{STMicroelectronics:2021b}

\bigskip

\subsection{Elektrische Eigenschaften}
Das MP34DT05 arbeitet mit einer Spannungsversorgung von 1,6 V bis 3,6 V und verbraucht im Normalbetrieb 650 µA. Im Energiesparmodus sinkt der Stromverbrauch auf $\leq 5\,\mu\text{A}$, was die Energieeffizienz für batteriebetriebene Geräte optimiert (Datenblatt MP34DT05, S. 6, Tab. 3). Die integrierte Analog-Signal-Wandlung ermöglicht eine direkte Ausgabe des Audiosignals im PDM-Format, das typischerweise als 1-Bit Datenstrom verarbeitet wird. Dies eliminiert die Notwendigkeit externer ADC-Komponenten und vereinfacht das Schaltungsdesign (Datenblatt MP34DT05, S. 1)\cite{STMicroelectronics:2021b}.

\bigskip

\subsection{Mechanisches Design und Umweltverträglichkeit}
Das HCLA-Gehäuse ist SMD-tauglich, EMI geschützt und verfügt über eine Top-Port-Design-Schallöffnung, die eine direkte Schallzufuhr von der Oberseite ermöglicht. Die Abmessungen des HLCA-Gehäuses (Länge: 3,0 mm ±0,1 mm, Breite: 4,0 mm ± 0,1 mm, Höhe: 1,0 mm ± 0,1 mm) ermöglichen eine einfache Integration in kompakte Systeme (Datenblatt MP34DT05, S. 18, Abb. 11). Das MP34DT05 entspricht den Umweltstandards ECOPACK®, RoHS und Green und ist somit frei von schädlichen Substanzen und recyclingfreundlich \cite{STMicroelectronics:2021b}.

\bigskip

\subsection{Funktionen und Anwendungen}
Das MP34DT05 unterstützt eine Stereo-Konfiguration durch die L/R- Kanalauswahl über Pin 2 (Datenblatt MP34DT05, S. 16, Tab. 7). Typische Anwendungen hierfür sind mobile Endgeräte, Laptops und Notebooks, Tragbare Mediaplayer, VoIP, Spracherkennung, A/V-E-Learning-Geräte oder Gaming- und Virtual-Reality-Eingabegeräte \cite{STMicroelectronics:2021b}.

\bigskip

\subsection{Betriebsbedingungen}
Das MP34DT05 ist für einen Temperaturbereich von -40 °C bis +85 °C ausgelegt und widersteht mechanischen Belastungen sowie elektrostatischen Entladungen (ESD-Schutz bis ± 2000 V HBM). Für die Montage wird ein SMD-kompatibles Lötprofil gemäß JEDEC J-STD-020 empfohlen \cite{STMicroelectronics:2021b}.


\section{Spezifikationen}

Im folgenden Abschnitt werden die technischen Spezifikationen des MP34DT05 aufgeführt.

\bigskip

\subsection{Allgemein}
\begin{itemize}
	\item Typ: Digitales MEMS-Mikrofon
	\item Richtcharakteristik: Omnidirektional
	\item Ausgangssignal: PDM (1-Bit-Stream)
	\item Betriebstemperatur: -40 °C bis +85 °C
\end{itemize}

\bigskip

\subsection{Elektrisch}
\begin{itemize}
	\item Spannungsversorgung: 1,6 V bis 3,6 V
	\item Stromverbrauch:
	\begin{itemize}
		\item Normalmodus: 650 µA
		\item Energiesparmodus: $\leq 5\,\mu\text{A}$
	\end{itemize}
	\item Integrierter ADC: Ja (direkte PDM-Digitalausgabe)
\end{itemize}

\subsection{Akustisch}
\begin{itemize}
	\item Frequenzbereich: 100 Hz bis 10 kHz
	\item Empfindlichkeit: -26 dBFS ± 3 dB (bei 94 dB SPL)
	\item Signal-Rausch-Abstand: 64 dB (A-gewichtet)
	\item Max. Schalldruck (AOP): 122,5 dB SPL
\end{itemize}

\subsection{Mechanisch}
\begin{itemize}
	\item Gehäuse: HCLGA (4\,mm × 3\,mm × 1\,mm)
	%\item Gewicht: Nicht angegeben
	\item Montage: SMD-tauglich, EMI-geschützt
\end{itemize}

\subsection{Funktionen}
\begin{itemize}
	\item Omnidirektionale Schallaufnahme
	\item L/R-Kanalauswahl für Stereo-Anwendungen
	\item Umweltzertifizierungen: RoHS, ECOPACK®, „Green“
\end{itemize}

\cite{STMicroelectronics:2021b}

\section{PDM Bibliothek} \label{sec:PDM.h}

\subsection{Beschreibung}

Die \texttt{PDM.h}-Bibliothek ermöglicht die direkte Ansteuerung des MP34DT05-MEMS(Micro Electron Mechanical System)-Mikrofons auf dem Arduino Nano 33 BLE Sense Lite. Sie nutzt die hardwarebeschleunigte PDM-Schnittstelle des nRF52840-Mikrocontrollers, um Pulse-Density-Modulation (PDM)-Daten in ein 16-Bit-PCM-Format umzuwandeln. Dies geschieht automatisch im Hintergrund, sodass Entwickler:innen sich auf die Implementierung von Anwendungen wie Spracherkennung, Audioanalyse oder maschinellem Lernen auf Edge-Geräten konzentrieren können \cite{ArduinoPDM:2023}.

Die Bibliothek ist für Echtzeitanwendungen optimiert und verwendet DMA (Direct Memory Access), um die CPU-Last zu minimieren. Dadurch bleibt genügend Rechenleistung für parallele Aufgaben wie Sensordatenverarbeitung oder Kommunikation mit anderen Geräten erhalten. Die Konfiguration der Abtastrate (zwischen 16 kHz und 48 kHz) sowie die Anpassung der Audiopuffer erfolgen über eine vereinfachte API, die speziell für die Integration in Arduino-Projekte entwickelt wurde \cite{ArduinoPDM:2023}.

Dank der engen Hardware-Integration läuft die Audiobearbeitung stabil und mit minimaler Latenz, was die Bibliothek ideal für IoT- und Edge-Computing-Szenarien macht. Die Umwandlung von PDM zu PCM erfolgt direkt auf Controller-Ebene, ohne zusätzliche externe Komponenten oder komplexe Schaltungen \cite{ArduinoPDM:2023}.

\subsection{Installation}

Die Installation der PDM-Bibliothek kann auf zwei Arten erfolgen:

\subsubsection{Über den Board Manager:}
	\begin{itemize}
		\item Öffne die Arduino IDE.
		\item Navigiere zu \textbf{Tools > Board > Boards Manager}.
		\item Suche nach \textit{Arduino mbed OS Giga Boards} und installiere die Version 1.2.3 von Adafruit.
		\item Suche im Bibliotheks-Manager nach \texttt{PDM} oder \texttt{Arduino Sound} und installiere die Bibliothek.
	\end{itemize}
	
\subsubsection{Über eine ZIP-Bibliothek:}
	\begin{itemize}
		\item Lade die \texttt{PDM}-Bibliothek von GitHub herunter \cite{Arduino:2025PDMlib}.
		\item In der Arduino IDE: \textbf{Sketch > Include Library > Add .ZIP Library…}
		\item Wähle die heruntergeladene ZIP-Datei aus.
	\end{itemize}

Nach erfolgreicher Installation kann die Funktionalität über \textbf{File > Examples > PDM} überprüft werden.


%\textcolor{red}{
%\begin{itemize}
%    \item data types
%    \item data structure
%    \item data quantity
%    \item data quality
%\end{itemize}
%}

\subsection{Funktionen}

Die Bibliothek bietet folgende Funktionen \cite{ArduinoPDM:2023}:

\begin{itemize}
	\item \PYTHON{PDM.begin(channels, sampleRate)}: Initialisiert die PDM-Schnittstelle mit der gewünschten Kanalanzahl (1 für Mono, 2 für Stereo) und Abtastrate (z. B. 16 kHz, 48 kHz).
	\item \PYTHON{PDM.end()}: Beendet die Audioaufnahme und gibt alle Hardware-Ressourcen des Mikrocontrollers frei.
	\item \PYTHON{PDM.available()}: Gibt die Anzahl der aktuell im Puffer verfügbaren Audiobytes zurück.
	\item \PYTHON{PDM.read(buffer, bufferSize)}: Liest PCM-Daten in einen nutzerdefinierten Puffer zur Weiterverarbeitung (z. B. für Spracherkennung oder Speicherung).
	\item \PYTHON{PDM.onReceive(callback)}: Registriert eine Callback-Funktion, die automatisch bei neuen Audio-Daten ausgelöst wird (Echtzeitverarbeitung).
	\item \PYTHON{PDM.setGain(wert)}: Passt die Mikrofonempfindlichkeit an (Standard: 20, Wertebereich 0–255). Unterstützt nicht alle Boards.
	\item \PYTHON{PDM.setBufferSize(size)}: Definiert die Größe des internen Audio-Puffers in Bytes (muss vor \PYTHON{PDM.begin()} gesetzt werden).
\end{itemize}



%\section{Example}

%In this section, a simple example is described in detail, which demonstrates how the library supports the sensor/actor.


%\subsection{Manual}

%\subsection{Inside the Example}

%\subsection{Code}

%\subsection{Files}



\section{Kalibrierung}

Der folgende Code \ref{lst:pdm_microphone_1} dient der Initialisierung eines PDM-Mikrofons eines Arduinos und der automatischen Kalibrierung eines Schwellenwerts (Threshold), der später beispielsweise zur Geräuscherkennung oder Sprachaktivierung genutzt werden kann. Der Schwellenwert ermöglicht es, zwischen Umgebungsgeräuschen und tatsächlichen, relevanten Geräuschereignissen zu unterscheiden.
Zu Beginn wird die benötigte Bibliothek \PYTHON{<PDM.h>} eingebunden, welche die Kommunikation mit dem Mikrofonmodul ermöglicht. Ein Puffer mit 256 Einträgen vom Typ \PYTHON{short} dient zur Zwischenspeicherung der eingehenden Audiodaten. Zusätzlich wird eine \PYTHON{volatile} Variable \PYTHON{samplesRead} deklariert, welche speichert, wie viele Samples aktuell aus dem Mikrofon gelesen wurden. Die Verwendung des \PYTHON{volatile}-Schlüsselworts stellt sicher, dass der Compiler berücksichtigt, dass sich dieser Wert jederzeit durch eine Interrupt-gesteuerte Callback-Funktion ändern kann.
Ein weiterer zentraler Bestandteil ist die Variable\PYTHON{ threshold}, in welcher der Schwellenwert abgelegt wird, der im Rahmen der Kalibrierung ermittelt wird. Die Kalibrierung erfolgt in einer separaten Funktion \PYTHON{calibrateThreshold()}, die über einen Zeitraum von zwei Sekunden die durchschnittliche Lautstärke der Umgebung misst. Dazu wird eine Summe der Absolutwerte aller eingehenden Audiosamples gebildet und anschließend durch die Anzahl der Samples geteilt. Um auf potenziell lautere Geräuschereignisse vorbereitet zu sein, wird der Mittelwert mit einem Faktor von drei multipliziert. Sollte während der Kalibrierungszeit kein einziger Wert erfasst werden, greift ein Standardwert von 15000 als Fallback.
Die Funktion \PYTHON{onPDMdata()} fungiert als Callback und wird immer dann automatisch aufgerufen, wenn neue Audiodaten vom Mikrofon zur Verfügung stehen. Sie liest 256 Samples aus dem Mikrofon in den definierten Sample-Puffer und aktualisiert die Anzahl der gelesenen Samples.
Im \PYTHON{setup()}-Block wird die serielle Verbindung initialisiert, um dem Benutzer Rückmeldungen über den Kalibrierungsprozess zu geben. Anschließend wird das Mikrofon gestartet und die Callback-Funktion mit dem PDM-Subsystem verknüpft. Gelingt die Initialisierung des Mikrofons nicht, bleibt das Programm in einer Endlosschleife stehen. Danach erfolgt die Kalibrierung des Schwellenwerts mittels der zuvor beschriebenen Methode. Der berechnete Threshold-Wert wird schließlich über die serielle Schnittstelle ausgegeben.
Die \PYTHON{loop()}-Funktion bleibt in diesem Beispiel leer, da der gesamte Fokus auf der Initialisierung und Kalibrierung liegt. In einer erweiterten Version dieses Programms könnte man hier beispielsweise eine Geräuscherkennung auf Basis des ermittelten Schwellenwerts implementieren, etwa um auf Klatschen, Sprachaktivierung oder andere akustische Ereignisse zu reagieren.
Insgesamt bildet dieser Code eine Grundlage für jedes Projekt, das eine Geräuscherkennung mithilfe eines PDM-Mikrofons erfordert. Durch die automatische Kalibrierung ist das System flexibel an unterschiedliche Umgebungsbedingungen anpassbar und somit robuster gegenüber Störgeräuschen.

{
	\captionof{code}{PDM Microphone Kalibrierung}\label{lst:pdm_microphone_1}
	\ArduinoExternal{}{../../Code/Arduino/MicrophoneMP34DT05/Calibration/Calibration.ino}
}

\section{Einfacher Code}

Im folgenden Code \ref{lst:pdm_microphone_2} wird die simple Initialisierung und kontinuierlichen Abtastung eines PDM-Mikrofons zur Erfassung von Audiosignalen. Zu Beginn werden einige Konstanten definiert: die Abtastrate (\PYTHON{SAMPLE RATE}) beträgt 16.000 Hz, es wird ein Mono-Kanal verwendet ((\PYTHON{CHANNELS = 1}), und der Pufferspeicher ist auf 256 16-Bit-Samples begrenzt.
Im Block \PYTHON{setup()} wird zunächst die serielle Kommunikation mit einer Baudrate von 115200 eingerichtet, um spätere Ausgaben über den seriellen Monitor zu ermöglichen. Anschließend wird das Mikrofon über die PDM-Bibliothek mit den zuvor definierten Parametern konfiguriert. Sollte die Initialisierung fehlschlagen, wird über eine separate Hilfsfunktion eine Fehlermeldung ausgegeben und das Programm gestoppt. Bei erfolgreicher Einrichtung wird die Verstärkung des Mikrofons auf einen festen Wert von 40 gesetzt. Abschließend gibt das Programm die Konfigurationswerte, also Samplingrate, Kanalanzahl und Puffergröße zur Kontrolle im seriellen Monitor aus.
In der Funktion \PYTHON{loop()} wird kontinuierlich geprüft, ob neue Audiodaten vom Mikrofon verfügbar sind. Sobald Daten vorhanden sind, werden sie in den globalen Sample-Puffer gelesen. Damit die Ausgabe nicht zu schnell erfolgt und den seriellen Monitor überlastet, wird der Inhalt des Puffers nur alle 100 Millisekunden ausgegeben. Dabei werden die aufgenommenen Rohwerte der Audiosamples einzeln ausgegeben, was zum Beispiel für eine erste Analyse des akustischen Signals oder zur Kalibrierung eines Schwellwerts genutzt werden kann.
Allgemein visualisiert der Code die grundlegende Funktionsweise eines PDM-Mikrofons, also das Erkennen eines Geräusches und das Ausgeben eines Feedbacks.

{
	\captionof{code}{PDM-Microphone -- einfacher Code}\label{lst:pdm_microphone_2}
	\ArduinoExternal{}{../../Code/Arduino/MicrophoneMP34DT05/Sample/Sample.ino}
}

\section{Einfache Anwendung}

Der folgende Code \ref{lst:pdm_microphone_3} ist eine simple Anwendung für eine Signalverarbeitung mit dem onboard-Mikrofon MP34DT05 des Arduino Nano 33 BLE Sense Lite. Diese erkennt Klatschen und lässt dann eine LED für eine bestimmte Zeit leuchten. Diese Anwendung basiert auf PDM (Pulse Density Modulation) und PCM (Pulse-Code-Modulation) diese werden verwendet, um Audiodaten aufzunehmen und zu verarbeiten um daraus ein Klatschen zu erkennen.
Beim Start des Programms wird in der Funktion \PYTHON{setup()} die serielle Kommunikation initialisiert und das Mikrofon mithilfe der PDM-Bibliothek gestartet. Danach folgt eine Kalibrierung des Schwellenwertes (Threshold), bei dem über zwei Sekunden hinweg der durchschnittliche Umgebungsgeräuschpegel gemessen wird. Der ermittelte Wert wird mit drei multipliziert, um eine Schwelle festzulegen, ab der ein Klatschen erkannt wird. Falls keine Daten gelesen werden, wird ein Standardwert als Schwellenwert gesetzt.
Im Hauptprogramm, also in der Funktion \PYTHON{loop()}, prüft das System kontinuierlich, ob neue Mikrofondaten vorhanden sind. Sobald dies der Fall ist, wird die maximale Amplitude aus dem Audiopuffer berechnet. Dieser Spitzenwert wird in einem Verlaufspuffer gespeichert, der drei Einträge umfasst. Dadurch können plötzliche Ausschläge besser beurteilt werden. Wenn der höchste aktuelle Wert diesen Schwellenwert überschreitet und seit der letzten Erkennung genug Zeit vergangen ist (Cooldown-Zeit), wird angenommen, dass ein Klatschen stattgefunden hat. In diesem Fall wird die LED eingeschaltet, und die Zeit der Erkennung wird gespeichert, um schnelle Mehrfachauslösungen zu vermeiden. Eine Meldung mit der gemessenen Stärke (dimensionsloser, amplitudenbasierter Rohwert) des Klatschens wird über die serielle Schnittstelle ausgegeben.
Parallel dazu wird ständig überprüft, ob die LED bereits eingeschaltet ist und ob die vorgesehene Leuchtdauer (fünf Sekunden) abgelaufen ist. Falls ja, wird die LED automatisch wieder ausgeschaltet, und es wird eine Meldung „LED off“ ausgegeben.
Die Funktion \PYTHON{onPDMdata()} wird automatisch aufgerufen, sobald neue Mikrofon-Daten bereitstehen. Sie füllt den Sample-Puffer mit 256 neuen Werten und gibt an, wie viele Samples tatsächlich gelesen wurden. Die Verarbeitung dieser Daten erfolgt dann, wie beschrieben, in der Hauptschleife.
Zusammengefasst erkennt der Arduino über das onboard-Mikrofon Klatschgeräusche, kalibriert sich dabei automatisch auf die Umgebungslautstärke und schaltet bei einem erkannten Klatschen für 5 Sekunden eine LED ein. Dabei wird sichergestellt, dass nicht mehrfach hintereinander ausgelöst wird und Störgeräusche ignoriert werden können.

{
	\captionof{code}{PDM Microphone Einfache Anwendung}\label{lst:pdm_microphone_3}
	\ArduinoExternal{}{../../Code/Arduino/MicrophoneMP34DT05/SimpleApp/SimpleApp.ino}
}
	
\section{Tests}

%\subsection{Einfache Funktions-Tests}

%\subsection{alle Funktionen testen}

\subsection{\PYTHON{PDM.begin()} – Initialisierung des PDM-Moduls}
Diese Funktion richtet zunächst die serielle Kommunikation mit 9600 Baud ein und wartet, bis der serielle Monitor erreichbar ist. Anschließend wird versucht, das PDM-Modul im Mono Modus bei einer Abtastrate von 16 kHz zu initialisieren. Der Test bestätigt, ob die Initialisierung erfolgreich war, und gibt dazu eine entsprechende Nachricht über die serielle Schnittstelle aus \ref{lst:pdm_microphone_4}.

{
	\captionof{code}{PDM.begin() Testcode}\label{lst:pdm_microphone_4}
	\ArduinoExternal{}{../../Code/Arduino/MicrophoneMP34DT05/Tests/PDMbegin/PDMbegin.ino}
}

\subsection{\PYTHON{PDM.end()} – Beenden und erneute Initialisierung des PDM-Moduls}
Hier wird das PDM-Modul zuerst gestartet (wieder im Mono Modus, 16 kHz) und anschließend wieder beendet. Nachdem das Modul ordnungsgemäß gestoppt wurde, erfolgt ein erneuter Versuch, das Modul zu initialisieren. Dadurch wird überprüft, ob ein korrektes Beenden und erneutes Starten möglich ist, was für eine stabile Audio-Erfassung wichtig ist \ref{lst:pdm_microphone_5}.

{
	\captionof{code}{PDM.end() Testcode}\label{lst:pdm_microphone_5}
	\ArduinoExternal{}{../../Code/Arduino/MicrophoneMP34DT05/Tests/PDMend/PDMend.ino}
}

\subsection{\PYTHON{PDM.available()} – Überwachung der verfügbaren Audiodaten}
In diesem Testcode wird das PDM-Modul initialisiert und der serielle Monitor gestartet. Anschließend wird kontinuierlich ausgegeben, wie viele Bytes an Audiodaten aktuell im internen Puffer verfügbar sind. Diese Funktion bietet so eine einfache Möglichkeit, den Datenfluss vom Mikrofoneingang zu beobachten \ref{lst:pdm_microphone_6}.

{
	\captionof{code}{PDM.available() Testcode}\label{lst:pdm_microphone_6}
	\ArduinoExternal{}{../../Code/Arduino/MicrophoneMP34DT05/Tests/PDMavailable/PDMavailable.ino}
}

\subsection{\PYTHON{PDM.read()} – Einlesen und Anzeigen von Audiodaten}
Dieser Code testet den eigentlichen Lesevorgang: Nach erfolgreicher Initialisierung des PDM-Moduls wird in regelmäßigen Abständen geprüft, ob neue Audiodaten verfügbar sind. Falls ja, werden diese in einen festgelegten Puffer gelesen und die ersten fünf Byte in hexadezimaler Darstellung über die serielle Schnittstelle ausgegeben. So wird geprüft, ob die Datenübertragung vom Mikrofon korrekt funktioniert \ref{lst:pdm_microphone_7}.

{
	\captionof{code}{PDM.read() Testcode}\label{lst:pdm_microphone_7}
	\ArduinoExternal{}{../../Code/Arduino/MicrophoneMP34DT05/Tests/PDMread/PDMread.ino}
}

\subsection{\PYTHON{PDM.onReceive()} – Interrupt-gesteuerte Datenübernahme}
Bei diesem Test steht das ereignisgesteuerte Einlesen von Audiodaten im Fokus. Ein Callback (onDataAvailable) wird registriert, das ausgelöst wird, sobald neue Daten im PDM-Puffer verfügbar sind. Im Anschluss wird in der Hauptschleife (loop) überprüft, ob das Callback aktiviert wurde, und ein entsprechender Hinweis über die serielle Verbindung ausgegeben. Damit wird die Reaktionsfähigkeit des Systems bei eingehenden Audiodaten getestet \ref{lst:pdm_microphone_8}.

{
	\captionof{code}{PDM.onRecieve() Testcode}\label{lst:pdm_microphone_8}
	\ArduinoExternal{}{../../Code/Arduino/MicrophoneMP34DT05/Tests/PDMonRecieve/PDMonRecieve.ino}
}

\subsection{\PYTHON{PDM.setGain()} – Konfiguration der Mikrofonempfindlichkeit}
Dieser Test zeigt, wie vor der Initialisierung des PDM Moduls der Mikrofon-Gain eingestellt werden kann. Der angegebene Wert (hier 255) soll die Empfindlichkeit des Mikrofons regulieren, sodass eine spätere Überprüfung der tatsächlichen Pegeländerung erfolgen kann. Die serielle Ausgabe informiert darüber, dass der Gain konfiguriert und das Modul erfolgreich gestartet wurde \ref{lst:pdm_microphone_9}.

{
	\captionof{code}{PDM.setGain() Testcode}\label{lst:pdm_microphone_9}
	\ArduinoExternal{}{../../Code/Arduino/MicrophoneMP34DT05/Tests/PDMsetGain/PDMsetGain.ino}
}

\subsection{\PYTHON{PDM.setBufferSize()} – Anpassung der Puffergröße}
Bei dieser Funktion wird vor dem Starten des PDM Moduls die Größe des internen Datenpuffers auf einen gewünschten Wert (z.B 1024 Byte) gesetzt. Dadurch kann getestet werden, ob die Puffergröße den Anforderungen angepasst werden kann, bevor das Modul in Betrieb genommen wird. Eine serielle Meldung bestätigt, dass die neue Puffergröße konfiguriert wurde, bevor das Modul initialisiert wird \ref{lst:pdm_microphone_10}.

{
	\captionof{code}{PDM.setBufferSize() Testcode}\label{lst:pdm_microphone_10}
	\ArduinoExternal{}{../../Code/Arduino/MicrophoneMP34DT05/Tests/PDMsetBuffersize/PDMsetBuffersize.ino}
}

\section{Weiterführende Literatur}

Die folgenden Literaturhinweise bieten vertiefende und ergänzende Einblicke in die theoretischen Grundlagen sowie praktische Anwendungen der Audio- und Elektroakustik. Die ausgewählten Werke ermöglichen eine weiterführende Beschäftigung mit einzelnen Aspekten der Audiotechnik und bieten umfassende Perspektiven, die über die Inhalte dieser Dokumentation hinausgehen.

\begin{itemize}
	\item Bernstein, Herbert: \textsl{Elektroakustik}. Wiesbaden:Springer Vieweg, 2019. \cite{Bernstein:2019d}
		
	Das Buch bietet vertiefende Einblicke in die physikalischen Grundlagen der Elektroakustik. Die umfassende Behandlung der theoretischen Zusammenhänge und Prinzipien der Schallwandlung ermöglicht eine weiterführende Auseinandersetzung mit Themen, die über den in dieser Dokumentation behandelten Rahmen hinausgehen.
	
	\item Görne, Thomas: \textsl{Tontechnik}. München:Hanser, 2014. \cite{Görne:2014e}
	
	Görne stellt praxisorientiert zahlreiche professionelle Anwendungsbeispiele und konkrete Umsetzungshinweise der Tontechnik vor. Diese ergänzen die in dieser Dokumentation dargestellten Arduino-basierten Audioanwendungen sinnvoll und erweitern das Verständnis für praktische Einsatzmöglichkeiten in einem breiteren Kontext.
	
	\item Weinzierl, Stefan: \textsl{Handbuch der Audiotechnik}. Springer 2008. \cite{Weinzierl:2008}
	
	Das Handbuch von Weinzierl vermittelt detailliertes Wissen zu zeitgemäßen audiotechnischen Konzepten und aktuellen Entwicklungen in der Audiotechnik. Die vertiefenden Inhalte und technischen Analysen bieten zusätzliche Perspektiven, um behandelte Themenbereiche umfassender zu erschließen und auf einem fortgeschrittenen Niveau weiter zu vertiefen.
	
\end{itemize} 
